\documentclass[11pt,a4paper]{article}

% ---- Packages ----
\usepackage[utf8]{inputenc}
\usepackage[T1]{fontenc}
\usepackage{amsmath,amssymb,amsthm}
\usepackage{mathtools}
\usepackage[margin=1in,headheight=14pt]{geometry}
\usepackage{hyperref}
\usepackage{enumitem}
\usepackage{xcolor}
\usepackage{tcolorbox}
\usepackage{fancyhdr}
\usepackage{booktabs}

% ---- Hyperref ----
\hypersetup{
    colorlinks=true,
    linkcolor=red,
    citecolor=blue,
    urlcolor=blue
}

% ---- Header ----
\pagestyle{fancy}
\fancyhf{}
\fancyhead[L]{\textit{Probability Review}}
\fancyhead[R]{\thepage}
\renewcommand{\headrulewidth}{0.4pt}

% ---- Theorem Environments ----
\theoremstyle{definition}
\newtheorem{definition}{Definition}[section]
\newtheorem{example}{Example}[section]
\newtheorem{exercise}{Exercise}[section]

\theoremstyle{plain}
\newtheorem{theorem}{Theorem}[section]
\newtheorem{lemma}[theorem]{Lemma}
\newtheorem{proposition}[theorem]{Proposition}
\newtheorem{corollary}[theorem]{Corollary}

\theoremstyle{remark}
\newtheorem{remark}{Remark}[section]

% ---- Colored Boxes ----
\tcbuselibrary{theorems,skins,breakable}

\newtcolorbox{keyresult}[1][]{
    colback=green!5!white,
    colframe=green!50!black,
    fonttitle=\bfseries,
    title=#1,
    breakable
}

\newtcolorbox{rlconnection}[1][]{
    colback=blue!5!white,
    colframe=blue!50!black,
    fonttitle=\bfseries,
    title=#1,
    breakable
}

\newtcolorbox{warning}[1][]{
    colback=red!5!white,
    colframe=red!50!black,
    fonttitle=\bfseries,
    title=#1,
    breakable
}

% ---- Operators ----
\newcommand{\R}{\mathbb{R}}
\newcommand{\N}{\mathbb{N}}
\newcommand{\Z}{\mathbb{Z}}
\newcommand{\E}{\mathbb{E}}
\newcommand{\Prob}{\mathbb{P}}
\DeclareMathOperator{\Var}{Var}
\DeclareMathOperator{\Cov}{Cov}
\newcommand{\indicator}{\mathbf{1}}

% ---- Title ----
\title{\textbf{Probability Review}\\[10pt]
\large Random Variables, Expectation, and Limit Theorems\\[5pt]
\normalsize Phase 0: Prerequisites Review $\mid$ Document 3 of 3\\
\normalsize Primary texts: Ross, \emph{A First Course in Probability}; Blitzstein \& Hwang, \emph{Introduction to Probability}}
\author{Curriculum: A Rigorous Learning Path to Multi-Agent Deep RL}
\date{February 2026}

\begin{document}

\maketitle

\begin{abstract}
This document is a self-contained refresher of classical (non-measure-theoretic) probability
theory. We cover sample spaces, conditional probability, random variables, expectation,
variance, joint distributions, moment generating functions, concentration inequalities, and
limit theorems. Every result is stated precisely and proved completely. The treatment is
elementary: we use probability mass functions, probability density functions, sums, and
integrals --- not $\sigma$-algebras or Lebesgue integration. Phase~01 will rebuild the entire
theory on measure-theoretic foundations; the goal here is to develop fluent computational
facility with the classical tools.
\end{abstract}

\tableofcontents
\newpage

%% ============================================================
%% SECTION 1: INTRODUCTION
%% ============================================================
\section{Introduction}
\label{sec:intro}

This document is a rigorous refresher of classical probability at the undergraduate level.
It is designed for readers who have seen this material before but need to recall the key
definitions, theorem statements, and proof techniques before proceeding to the
measure-theoretic treatment in Phase~01.

\medskip\noindent\textbf{What this document covers.}
\begin{enumerate}[nosep]
    \item \textbf{Sample spaces and events} --- axioms of probability, inclusion-exclusion,
          continuity of probability.
    \item \textbf{Conditional probability and independence} --- Bayes' theorem, mutual
          independence.
    \item \textbf{Discrete random variables} --- PMFs, Bernoulli, Binomial, Geometric, Poisson.
    \item \textbf{Continuous random variables} --- PDFs, Uniform, Exponential, Normal.
    \item \textbf{Expectation} --- linearity, LOTUS, indicator random variables.
    \item \textbf{Variance and covariance} --- computational formulas, variance of sums.
    \item \textbf{Joint distributions} --- marginals, conditional distributions, laws of total
          expectation and total variance.
    \item \textbf{Moment generating functions} --- computing moments, sums of independent
          random variables.
    \item \textbf{Probability inequalities} --- Markov, Chebyshev, Cauchy--Schwarz, Jensen.
    \item \textbf{Limit theorems} --- weak law of large numbers, central limit theorem.
    \item \textbf{Tail bounds and confidence intervals} --- Hoeffding's inequality, sample
          complexity.
\end{enumerate}

\medskip\noindent\textbf{What this document does not cover.}
We do not use $\sigma$-algebras, measurability, Lebesgue integration, filtrations,
or martingale theory. These are developed from scratch in Phase~01.

\begin{rlconnection}[Connection to Reinforcement Learning]
Reinforcement learning is fundamentally about decision-making under uncertainty. Every
concept here --- conditional probability, expectation, variance, the law of large
numbers --- has a direct RL counterpart. Rewards are random variables. Value functions are
conditional expectations. Monte Carlo methods converge by the law of large numbers. This
review ensures you can compute fluently before Phase~01 introduces the measure-theoretic
machinery that makes everything rigorous.
\end{rlconnection}

%% ============================================================
%% SECTION 2: SAMPLE SPACES AND EVENTS
%% ============================================================
\section{Sample Spaces and Events}
\label{sec:sample_spaces}

\begin{definition}[Sample Space, Outcome, Event]
\label{def:sample_space}
A \emph{sample space} $\Omega$ is a non-empty set whose elements $\omega \in \Omega$ are
called \emph{outcomes}. An \emph{event} is a subset $A \subseteq \Omega$.
\end{definition}

\begin{remark}
At this elementary level, we take the collection of all events to be the power set
$2^\Omega$ (all subsets of $\Omega$). In Phase~01, this will be replaced by a
$\sigma$-algebra $\mathcal{F} \subseteq 2^\Omega$, which is necessary for uncountable
sample spaces.
\end{remark}

\begin{definition}[Set Operations on Events]
\label{def:set_ops}
Let $A, B \subseteq \Omega$.
\begin{enumerate}[nosep]
    \item \emph{Union:} $A \cup B = \{\omega \in \Omega : \omega \in A \text{ or } \omega \in B\}$ (``$A$ or $B$ occurs'').
    \item \emph{Intersection:} $A \cap B = \{\omega \in \Omega : \omega \in A \text{ and } \omega \in B\}$ (``$A$ and $B$ occur'').
    \item \emph{Complement:} $A^c = \Omega \setminus A = \{\omega \in \Omega : \omega \notin A\}$ (``$A$ does not occur'').
    \item \emph{Disjoint:} $A$ and $B$ are \emph{disjoint} (or \emph{mutually exclusive}) if $A \cap B = \emptyset$.
\end{enumerate}
\end{definition}

\subsection{Axioms of Probability}

\begin{definition}[Probability (Kolmogorov Axioms)]
\label{def:probability}
A \emph{probability} on a sample space $\Omega$ is a function $\Prob$ assigning to each
event $A \subseteq \Omega$ a real number $\Prob(A)$ satisfying:
\begin{enumerate}[nosep]
    \item \textbf{Non-negativity:} $\Prob(A) \geq 0$ for every event $A$.
    \item \textbf{Normalization:} $\Prob(\Omega) = 1$.
    \item \textbf{Countable additivity:} If $A_1, A_2, \ldots$ are pairwise disjoint events, then
    \[
        \Prob\!\left(\bigcup_{n=1}^\infty A_n\right) = \sum_{n=1}^\infty \Prob(A_n).
    \]
\end{enumerate}
\end{definition}

\subsection{Basic Properties}

\begin{proposition}[Complement Rule]
\label{prop:complement}
For any event $A$, $\Prob(A^c) = 1 - \Prob(A)$.
\end{proposition}

\begin{proof}
Since $A$ and $A^c$ are disjoint and $A \cup A^c = \Omega$, Axiom~3 gives
$\Prob(\Omega) = \Prob(A) + \Prob(A^c)$. By Axiom~2, $\Prob(\Omega) = 1$, so
$\Prob(A^c) = 1 - \Prob(A)$.
\end{proof}

\begin{corollary}
$\Prob(\emptyset) = 0$.
\end{corollary}

\begin{proof}
$\emptyset = \Omega^c$, so $\Prob(\emptyset) = 1 - \Prob(\Omega) = 1 - 1 = 0$.
\end{proof}

\begin{proposition}[Monotonicity]
\label{prop:monotonicity}
If $A \subseteq B$, then $\Prob(A) \leq \Prob(B)$.
\end{proposition}

\begin{proof}
Write $B = A \cup (B \setminus A)$, where $A$ and $B \setminus A$ are disjoint. By
Axiom~3, $\Prob(B) = \Prob(A) + \Prob(B \setminus A) \geq \Prob(A)$, since
$\Prob(B \setminus A) \geq 0$ by Axiom~1.
\end{proof}

\begin{theorem}[Inclusion-Exclusion for Two Events]
\label{thm:incl_excl_two}
For any events $A$ and $B$,
\[
    \Prob(A \cup B) = \Prob(A) + \Prob(B) - \Prob(A \cap B).
\]
\end{theorem}

\begin{proof}
Write $A \cup B = A \cup (B \setminus A)$, where $A$ and $B \setminus A$ are disjoint.
By Axiom~3, $\Prob(A \cup B) = \Prob(A) + \Prob(B \setminus A)$. Now write
$B = (A \cap B) \cup (B \setminus A)$, where $A \cap B$ and $B \setminus A$ are disjoint.
By Axiom~3, $\Prob(B) = \Prob(A \cap B) + \Prob(B \setminus A)$, so
$\Prob(B \setminus A) = \Prob(B) - \Prob(A \cap B)$. Substituting,
$\Prob(A \cup B) = \Prob(A) + \Prob(B) - \Prob(A \cap B)$.
\end{proof}

\begin{theorem}[Inclusion-Exclusion for $n$ Events]
\label{thm:incl_excl_n}
For events $A_1, \ldots, A_n$,
\[
    \Prob\!\left(\bigcup_{i=1}^n A_i\right) = \sum_{k=1}^n (-1)^{k+1}
    \sum_{1 \leq i_1 < \cdots < i_k \leq n} \Prob(A_{i_1} \cap \cdots \cap A_{i_k}).
\]
\end{theorem}

\begin{proof}
We proceed by induction on $n$.

\emph{Base case ($n = 1$):} $\Prob(A_1) = \Prob(A_1)$. Trivially holds.

\emph{Base case ($n = 2$):} This is Theorem~\ref{thm:incl_excl_two}.

\emph{Inductive step:} Assume the formula holds for $n-1$ events. Write
$\bigcup_{i=1}^n A_i = \left(\bigcup_{i=1}^{n-1} A_i\right) \cup A_n$. By
Theorem~\ref{thm:incl_excl_two},
\[
    \Prob\!\left(\bigcup_{i=1}^n A_i\right) = \Prob\!\left(\bigcup_{i=1}^{n-1} A_i\right)
    + \Prob(A_n) - \Prob\!\left(\left(\bigcup_{i=1}^{n-1} A_i\right) \cap A_n\right).
\]
Since $\left(\bigcup_{i=1}^{n-1} A_i\right) \cap A_n = \bigcup_{i=1}^{n-1}(A_i \cap A_n)$,
we can apply the inductive hypothesis to both $\Prob\!\left(\bigcup_{i=1}^{n-1} A_i\right)$
and $\Prob\!\left(\bigcup_{i=1}^{n-1}(A_i \cap A_n)\right)$. Applying the inductive hypothesis
to the first term gives the alternating sum over subsets of $\{1, \ldots, n-1\}$. Adding
$\Prob(A_n)$ accounts for the singleton $\{n\}$. Subtracting the inductive hypothesis
applied to $\bigcup_{i=1}^{n-1}(A_i \cap A_n)$ contributes exactly the terms involving
$A_n$ in subsets of size $\geq 2$, each with the correct sign $(-1)^{k+1}$. Collecting
all terms yields the formula for $n$ events.
\end{proof}

\subsection{Continuity of Probability}

\begin{theorem}[Continuity of Probability]
\label{thm:continuity_prob}
Let $A_1 \subseteq A_2 \subseteq A_3 \subseteq \cdots$ be an increasing sequence of
events. Then
\[
    \Prob\!\left(\bigcup_{n=1}^\infty A_n\right) = \lim_{n \to \infty} \Prob(A_n).
\]
\end{theorem}

\begin{proof}
Define $B_1 = A_1$ and $B_n = A_n \setminus A_{n-1}$ for $n \geq 2$. The events
$B_1, B_2, \ldots$ are pairwise disjoint (since $A_n \setminus A_{n-1}$ contains only
elements in $A_n$ but not in $A_{n-1}$, and these sets are nested), and
$\bigcup_{n=1}^\infty B_n = \bigcup_{n=1}^\infty A_n$. By countable additivity,
\[
    \Prob\!\left(\bigcup_{n=1}^\infty A_n\right) = \sum_{n=1}^\infty \Prob(B_n).
\]
Now, $\sum_{k=1}^n \Prob(B_k) = \Prob\!\left(\bigcup_{k=1}^n B_k\right) = \Prob(A_n)$,
where the first equality uses finite additivity (the $B_k$ are disjoint) and the second
uses $\bigcup_{k=1}^n B_k = A_n$. Therefore
\[
    \Prob\!\left(\bigcup_{n=1}^\infty A_n\right) = \sum_{n=1}^\infty \Prob(B_n)
    = \lim_{n \to \infty} \sum_{k=1}^n \Prob(B_k) = \lim_{n \to \infty} \Prob(A_n).
    \qedhere
\]
\end{proof}

%% ============================================================
%% SECTION 3: CONDITIONAL PROBABILITY AND INDEPENDENCE
%% ============================================================
\section{Conditional Probability and Independence}
\label{sec:cond_prob}

\subsection{Conditional Probability}

\begin{definition}[Conditional Probability]
\label{def:cond_prob}
If $\Prob(B) > 0$, the \emph{conditional probability of $A$ given $B$} is
\[
    \Prob(A \mid B) = \frac{\Prob(A \cap B)}{\Prob(B)}.
\]
\end{definition}

\begin{theorem}[Multiplication Rule]
\label{thm:mult_rule}
For events $A_1, \ldots, A_n$ with $\Prob(A_1 \cap \cdots \cap A_{n-1}) > 0$,
\[
    \Prob(A_1 \cap \cdots \cap A_n) = \Prob(A_1)\,\Prob(A_2 \mid A_1)\,
    \Prob(A_3 \mid A_1 \cap A_2) \cdots \Prob(A_n \mid A_1 \cap \cdots \cap A_{n-1}).
\]
\end{theorem}

\begin{proof}
By induction on $n$. For $n = 2$, the definition of conditional probability gives
$\Prob(A_1 \cap A_2) = \Prob(A_1)\,\Prob(A_2 \mid A_1)$.

Assume the result holds for $n - 1$ events. Then
\begin{align*}
    \Prob(A_1 \cap \cdots \cap A_n)
    &= \Prob(A_1 \cap \cdots \cap A_{n-1}) \cdot \Prob(A_n \mid A_1 \cap \cdots \cap A_{n-1}) \\
    &= \Prob(A_1)\,\Prob(A_2 \mid A_1) \cdots \Prob(A_{n-1} \mid A_1 \cap \cdots \cap A_{n-2})
       \cdot \Prob(A_n \mid A_1 \cap \cdots \cap A_{n-1}),
\end{align*}
where the second line applies the inductive hypothesis.
\end{proof}

\begin{theorem}[Law of Total Probability]
\label{thm:total_prob}
Let $B_1, B_2, \ldots, B_n$ be a partition of $\Omega$ (i.e., the $B_i$ are pairwise
disjoint and $\bigcup_{i=1}^n B_i = \Omega$) with $\Prob(B_i) > 0$ for all $i$. Then for
any event $A$,
\[
    \Prob(A) = \sum_{i=1}^n \Prob(A \mid B_i)\,\Prob(B_i).
\]
\end{theorem}

\begin{proof}
Since the $B_i$ partition $\Omega$, we have $A = A \cap \Omega = A \cap
\left(\bigcup_{i=1}^n B_i\right) = \bigcup_{i=1}^n (A \cap B_i)$. The events
$A \cap B_1, \ldots, A \cap B_n$ are pairwise disjoint (since the $B_i$ are). By
finite additivity,
\[
    \Prob(A) = \sum_{i=1}^n \Prob(A \cap B_i) = \sum_{i=1}^n \Prob(A \mid B_i)\,\Prob(B_i),
\]
where the last step uses $\Prob(A \cap B_i) = \Prob(A \mid B_i)\,\Prob(B_i)$.
\end{proof}

\begin{theorem}[Bayes' Theorem]
\label{thm:bayes}
Let $B_1, \ldots, B_n$ be a partition of $\Omega$ with $\Prob(B_i) > 0$ for all $i$, and
let $A$ be an event with $\Prob(A) > 0$. Then
\[
    \Prob(B_j \mid A) = \frac{\Prob(A \mid B_j)\,\Prob(B_j)}{\sum_{i=1}^n \Prob(A \mid B_i)\,\Prob(B_i)}.
\]
\end{theorem}

\begin{proof}
By definition, $\Prob(B_j \mid A) = \Prob(A \cap B_j)/\Prob(A)$. The numerator equals
$\Prob(A \mid B_j)\,\Prob(B_j)$ by the definition of conditional probability. The
denominator equals $\sum_{i=1}^n \Prob(A \mid B_i)\,\Prob(B_i)$ by the law of total
probability (Theorem~\ref{thm:total_prob}).
\end{proof}

\subsection{Independence}

\begin{definition}[Independence of Two Events]
\label{def:indep_two}
Events $A$ and $B$ are \emph{independent} if
\[
    \Prob(A \cap B) = \Prob(A)\,\Prob(B).
\]
\end{definition}

\begin{remark}
When $\Prob(B) > 0$, independence is equivalent to $\Prob(A \mid B) = \Prob(A)$: knowing
that $B$ occurred does not change the probability of $A$.
\end{remark}

\begin{definition}[Mutual Independence]
\label{def:mutual_indep}
Events $A_1, \ldots, A_n$ are \emph{mutually independent} if for every subset
$\{i_1, \ldots, i_k\} \subseteq \{1, \ldots, n\}$ with $k \geq 2$,
\[
    \Prob(A_{i_1} \cap \cdots \cap A_{i_k}) = \Prob(A_{i_1}) \cdots \Prob(A_{i_k}).
\]
This requires the product formula to hold for \emph{all} $2^n - n - 1$ subsets of size
$\geq 2$, not merely for the $\binom{n}{2}$ pairs.
\end{definition}

\begin{warning}[Pairwise $\neq$ Mutual Independence]
Pairwise independence does not imply mutual independence. The following example shows that
the product rule can hold for every pair yet fail for the triple.
\end{warning}

\begin{example}[Pairwise but Not Mutually Independent]
\label{ex:pairwise_not_mutual}
Toss two fair coins. Let $A = \{\text{first coin is heads}\}$,
$B = \{\text{second coin is heads}\}$, and $C = \{\text{exactly one head}\}$.
The sample space is $\{HH, HT, TH, TT\}$, each with probability $1/4$.

We have $\Prob(A) = \Prob(B) = \Prob(C) = 1/2$.
\begin{itemize}[nosep]
    \item $A \cap B = \{HH\}$, so $\Prob(A \cap B) = 1/4 = \Prob(A)\Prob(B)$.
    \item $A \cap C = \{HT\}$, so $\Prob(A \cap C) = 1/4 = \Prob(A)\Prob(C)$.
    \item $B \cap C = \{TH\}$, so $\Prob(B \cap C) = 1/4 = \Prob(B)\Prob(C)$.
\end{itemize}
All three pairs are independent. However, $A \cap B \cap C = \emptyset$, so
$\Prob(A \cap B \cap C) = 0 \neq 1/8 = \Prob(A)\Prob(B)\Prob(C)$.
Hence $A, B, C$ are pairwise independent but \emph{not} mutually independent.
\end{example}

\begin{rlconnection}[Connection to RL]
The Markov property --- the foundation of MDPs --- is an independence statement: the future
is independent of the past given the present. Bayes' theorem is the update rule in Bayesian
reinforcement learning and POMDP belief updates.
\end{rlconnection}

%% ============================================================
%% SECTION 4: DISCRETE RANDOM VARIABLES
%% ============================================================
\section{Discrete Random Variables}
\label{sec:discrete_rv}

\begin{definition}[Random Variable]
\label{def:rv}
A \emph{random variable} is a function $X \colon \Omega \to \R$ that assigns a real number
to each outcome. We write $\{X = x\}$ as shorthand for $\{\omega \in \Omega : X(\omega) = x\}$.
\end{definition}

\begin{remark}
At this level, the definition is purely informal: any function from $\Omega$ to $\R$.
Phase~01 will add the measurability requirement that makes this rigorous for uncountable
sample spaces.
\end{remark}

\begin{definition}[Probability Mass Function]
\label{def:pmf}
A random variable $X$ is \emph{discrete} if it takes values in a finite or countable set
$\{x_1, x_2, \ldots\}$. Its \emph{probability mass function (PMF)} is
\[
    p_X(x) = \Prob(X = x).
\]
The PMF satisfies $p_X(x) \geq 0$ for all $x$ and $\sum_x p_X(x) = 1$.
\end{definition}

\begin{definition}[Cumulative Distribution Function]
\label{def:cdf}
The \emph{cumulative distribution function (CDF)} of a random variable $X$ is
\[
    F_X(x) = \Prob(X \leq x), \quad x \in \R.
\]
For a discrete random variable, $F_X(x) = \sum_{x_i \leq x} p_X(x_i)$.
\end{definition}

\subsection{Common Discrete Distributions}

\begin{definition}[Bernoulli Distribution]
\label{def:bernoulli}
A random variable $X$ has the $\mathrm{Bernoulli}(p)$ distribution, $0 \leq p \leq 1$, if
$X$ takes values in $\{0, 1\}$ with
\[
    p_X(1) = p, \quad p_X(0) = 1 - p.
\]
\end{definition}

\begin{proposition}
\label{prop:bernoulli_moments}
If $X \sim \mathrm{Bernoulli}(p)$, then $\E[X] = p$ and $\Var(X) = p(1-p)$.
\end{proposition}

\begin{proof}
$\E[X] = 0 \cdot (1-p) + 1 \cdot p = p$.
For the variance: $\E[X^2] = 0^2 \cdot (1-p) + 1^2 \cdot p = p$, so
$\Var(X) = \E[X^2] - (\E[X])^2 = p - p^2 = p(1-p)$.
\end{proof}

\begin{definition}[Binomial Distribution]
\label{def:binomial}
Suppose we perform $n$ independent Bernoulli trials, each with success probability $p$.
Let $X$ count the total number of successes. Then $X \sim \mathrm{Binomial}(n, p)$.
\end{definition}

\begin{proposition}[Binomial PMF]
\label{prop:binomial_pmf}
If $X \sim \mathrm{Binomial}(n, p)$, then
\[
    p_X(k) = \binom{n}{k} p^k (1-p)^{n-k}, \quad k = 0, 1, \ldots, n.
\]
\end{proposition}

\begin{proof}
Each specific sequence of $k$ successes and $n - k$ failures has probability
$p^k(1-p)^{n-k}$ by independence. There are $\binom{n}{k}$ such sequences (choosing which
$k$ of the $n$ trials are successes), so $\Prob(X = k) = \binom{n}{k} p^k(1-p)^{n-k}$.
\end{proof}

\begin{proposition}
\label{prop:binomial_moments}
If $X \sim \mathrm{Binomial}(n, p)$, then $\E[X] = np$ and $\Var(X) = np(1-p)$.
\end{proposition}

\begin{proof}
Write $X = X_1 + \cdots + X_n$, where $X_i \sim \mathrm{Bernoulli}(p)$ independently. By
linearity of expectation (Theorem~\ref{thm:linearity} below),
$\E[X] = \sum_{i=1}^n \E[X_i] = np$. Since the $X_i$ are independent, $\Var(X) =
\sum_{i=1}^n \Var(X_i) = np(1-p)$ (using Corollary~\ref{cor:var_indep_sum} below).
\end{proof}

\begin{definition}[Geometric Distribution]
\label{def:geometric}
Perform independent Bernoulli trials with success probability $p > 0$. Let $X$ be the trial
number of the first success. Then $X \sim \mathrm{Geometric}(p)$.
\end{definition}

\begin{proposition}[Geometric PMF]
\label{prop:geometric_pmf}
If $X \sim \mathrm{Geometric}(p)$, then
\[
    p_X(k) = (1-p)^{k-1} p, \quad k = 1, 2, 3, \ldots
\]
\end{proposition}

\begin{proof}
The first success occurs on trial $k$ exactly when the first $k-1$ trials are failures and
trial $k$ is a success. By independence, this probability is $(1-p)^{k-1} \cdot p$.
\end{proof}

\begin{proposition}
\label{prop:geometric_moments}
If $X \sim \mathrm{Geometric}(p)$, then $\E[X] = 1/p$ and $\Var(X) = (1-p)/p^2$.
\end{proposition}

\begin{proof}
Let $q = 1 - p$. We compute $\E[X]$ directly:
\[
    \E[X] = \sum_{k=1}^\infty k \, q^{k-1} p = p \sum_{k=1}^\infty k \, q^{k-1}.
\]
Recall that $\sum_{k=1}^\infty k \, q^{k-1} = \frac{d}{dq}\sum_{k=0}^\infty q^k
= \frac{d}{dq}\frac{1}{1-q} = \frac{1}{(1-q)^2} = \frac{1}{p^2}$ for $|q| < 1$.
Thus $\E[X] = p/p^2 = 1/p$.

For the variance, compute $\E[X^2]$. We use $\E[X(X-1)] = \sum_{k=1}^\infty k(k-1)q^{k-1}p
= p \, q \sum_{k=2}^\infty k(k-1)q^{k-2}$. Since $\sum_{k=2}^\infty k(k-1)q^{k-2}
= \frac{d^2}{dq^2}\frac{1}{1-q} = \frac{2}{(1-q)^3} = \frac{2}{p^3}$, we get
$\E[X(X-1)] = 2q/p^2$. Then $\E[X^2] = \E[X(X-1)] + \E[X] = 2q/p^2 + 1/p$, and
\[
    \Var(X) = \E[X^2] - (\E[X])^2 = \frac{2q}{p^2} + \frac{1}{p} - \frac{1}{p^2}
    = \frac{2q + p - 1}{p^2} = \frac{q}{p^2} = \frac{1-p}{p^2},
\]
where we used $2q + p - 1 = 2(1-p) + p - 1 = 1 - p = q$.
\end{proof}

\begin{definition}[Poisson Distribution]
\label{def:poisson}
A random variable $X$ has the $\mathrm{Poisson}(\lambda)$ distribution, $\lambda > 0$, if
\[
    p_X(k) = \frac{e^{-\lambda}\lambda^k}{k!}, \quad k = 0, 1, 2, \ldots
\]
\end{definition}

\begin{theorem}[Poisson as a Limit of Binomials]
\label{thm:poisson_limit}
Let $X_n \sim \mathrm{Binomial}(n, \lambda/n)$ with $\lambda > 0$ fixed. Then for each
$k = 0, 1, 2, \ldots$,
\[
    \lim_{n \to \infty} \Prob(X_n = k) = \frac{e^{-\lambda}\lambda^k}{k!}.
\]
\end{theorem}

\begin{proof}
With $p = \lambda/n$,
\[
    \Prob(X_n = k) = \binom{n}{k}\left(\frac{\lambda}{n}\right)^k
    \left(1 - \frac{\lambda}{n}\right)^{n-k}.
\]
Write $\binom{n}{k} = \frac{n(n-1)\cdots(n-k+1)}{k!}$, so
\[
    \Prob(X_n = k) = \frac{n(n-1)\cdots(n-k+1)}{n^k} \cdot \frac{\lambda^k}{k!}
    \cdot \left(1 - \frac{\lambda}{n}\right)^{n-k}.
\]
As $n \to \infty$: $\frac{n(n-1)\cdots(n-k+1)}{n^k} \to 1$ (each factor $\frac{n-j}{n}
\to 1$), and $\left(1 - \frac{\lambda}{n}\right)^n \to e^{-\lambda}$, and
$\left(1 - \frac{\lambda}{n}\right)^{-k} \to 1$. Therefore $\Prob(X_n = k) \to
e^{-\lambda}\lambda^k/k!$.
\end{proof}

\begin{proposition}
\label{prop:poisson_moments}
If $X \sim \mathrm{Poisson}(\lambda)$, then $\E[X] = \lambda$ and $\Var(X) = \lambda$.
\end{proposition}

\begin{proof}
\begin{align*}
    \E[X] &= \sum_{k=0}^\infty k \cdot \frac{e^{-\lambda}\lambda^k}{k!}
    = \sum_{k=1}^\infty \frac{e^{-\lambda}\lambda^k}{(k-1)!}
    = \lambda e^{-\lambda} \sum_{j=0}^\infty \frac{\lambda^j}{j!}
    = \lambda e^{-\lambda} \cdot e^{\lambda} = \lambda.
\end{align*}
For the variance, compute $\E[X(X-1)]$:
\begin{align*}
    \E[X(X-1)] &= \sum_{k=2}^\infty k(k-1)\frac{e^{-\lambda}\lambda^k}{k!}
    = \lambda^2 e^{-\lambda} \sum_{j=0}^\infty \frac{\lambda^j}{j!} = \lambda^2.
\end{align*}
Then $\E[X^2] = \E[X(X-1)] + \E[X] = \lambda^2 + \lambda$, so
$\Var(X) = \lambda^2 + \lambda - \lambda^2 = \lambda$.
\end{proof}

%% ============================================================
%% SECTION 5: CONTINUOUS RANDOM VARIABLES
%% ============================================================
\section{Continuous Random Variables}
\label{sec:continuous_rv}

\begin{definition}[Probability Density Function]
\label{def:pdf}
A random variable $X$ is \emph{continuous} if there exists a non-negative function
$f_X \colon \R \to [0, \infty)$ such that for every interval $[a, b]$,
\[
    \Prob(a \leq X \leq b) = \int_a^b f_X(x)\,dx.
\]
The function $f_X$ is called the \emph{probability density function (PDF)} of $X$. The PDF
satisfies $\int_{-\infty}^\infty f_X(x)\,dx = 1$.
\end{definition}

\begin{proposition}
\label{prop:cdf_pdf}
If $X$ is continuous with PDF $f_X$, then its CDF is $F_X(x) = \int_{-\infty}^x f_X(t)\,dt$,
and $f_X(x) = F_X'(x)$ wherever $F_X$ is differentiable.
\end{proposition}

\begin{proof}
By definition, $F_X(x) = \Prob(X \leq x) = \int_{-\infty}^x f_X(t)\,dt$. By the
fundamental theorem of calculus, $F_X'(x) = f_X(x)$ at every point where $f_X$ is
continuous.
\end{proof}

\begin{proposition}
\label{prop:continuous_point_zero}
If $X$ is a continuous random variable, then $\Prob(X = x) = 0$ for every $x \in \R$.
\end{proposition}

\begin{proof}
$\Prob(X = x) = \Prob(x \leq X \leq x) = \int_x^x f_X(t)\,dt = 0$.
\end{proof}

\subsection{Common Continuous Distributions}

\begin{definition}[Uniform Distribution]
\label{def:uniform}
$X \sim \mathrm{Uniform}(a, b)$ if $X$ has PDF
\[
    f_X(x) = \begin{cases} \frac{1}{b - a} & \text{if } a \leq x \leq b, \\ 0 & \text{otherwise.} \end{cases}
\]
\end{definition}

\begin{proposition}
\label{prop:uniform_moments}
If $X \sim \mathrm{Uniform}(a, b)$, then $\E[X] = \frac{a+b}{2}$ and
$\Var(X) = \frac{(b-a)^2}{12}$.
\end{proposition}

\begin{proof}
\[
    \E[X] = \int_a^b x \cdot \frac{1}{b-a}\,dx = \frac{1}{b-a}\cdot\frac{x^2}{2}\bigg|_a^b
    = \frac{b^2 - a^2}{2(b-a)} = \frac{a + b}{2}.
\]
\[
    \E[X^2] = \int_a^b x^2 \cdot \frac{1}{b-a}\,dx = \frac{1}{b-a}\cdot\frac{x^3}{3}\bigg|_a^b
    = \frac{b^3 - a^3}{3(b-a)} = \frac{a^2 + ab + b^2}{3}.
\]
Thus
\[
    \Var(X) = \E[X^2] - (\E[X])^2 = \frac{a^2 + ab + b^2}{3} - \frac{(a+b)^2}{4}
    = \frac{4(a^2 + ab + b^2) - 3(a+b)^2}{12} = \frac{a^2 - 2ab + b^2}{12}
    = \frac{(b-a)^2}{12}.
    \qedhere
\]
\end{proof}

\begin{definition}[Exponential Distribution]
\label{def:exponential}
$X \sim \mathrm{Exponential}(\lambda)$, $\lambda > 0$, if $X$ has PDF
\[
    f_X(x) = \begin{cases} \lambda e^{-\lambda x} & \text{if } x \geq 0, \\ 0 & \text{if } x < 0. \end{cases}
\]
\end{definition}

\begin{proposition}
\label{prop:exponential_moments}
If $X \sim \mathrm{Exponential}(\lambda)$, then $\E[X] = 1/\lambda$ and
$\Var(X) = 1/\lambda^2$.
\end{proposition}

\begin{proof}
Using integration by parts with $u = x$, $dv = \lambda e^{-\lambda x}\,dx$:
\[
    \E[X] = \int_0^\infty x \lambda e^{-\lambda x}\,dx
    = \left[-x e^{-\lambda x}\right]_0^\infty + \int_0^\infty e^{-\lambda x}\,dx
    = 0 + \frac{1}{\lambda} = \frac{1}{\lambda}.
\]
Similarly, integrating by parts twice:
\[
    \E[X^2] = \int_0^\infty x^2 \lambda e^{-\lambda x}\,dx = \frac{2}{\lambda^2}.
\]
Thus $\Var(X) = 2/\lambda^2 - 1/\lambda^2 = 1/\lambda^2$.
\end{proof}

\begin{theorem}[Memoryless Property of the Exponential]
\label{thm:memoryless}
If $X \sim \mathrm{Exponential}(\lambda)$, then for all $s, t \geq 0$,
\[
    \Prob(X > s + t \mid X > s) = \Prob(X > t).
\]
\end{theorem}

\begin{proof}
The survival function is $\Prob(X > x) = e^{-\lambda x}$ for $x \geq 0$. By definition
of conditional probability,
\[
    \Prob(X > s + t \mid X > s) = \frac{\Prob(X > s + t \text{ and } X > s)}{\Prob(X > s)}
    = \frac{\Prob(X > s + t)}{\Prob(X > s)} = \frac{e^{-\lambda(s+t)}}{e^{-\lambda s}}
    = e^{-\lambda t} = \Prob(X > t).
    \qedhere
\]
\end{proof}

\begin{definition}[Normal (Gaussian) Distribution]
\label{def:normal}
$X \sim \mathcal{N}(\mu, \sigma^2)$ if $X$ has PDF
\[
    f_X(x) = \frac{1}{\sigma\sqrt{2\pi}} \exp\!\left(-\frac{(x - \mu)^2}{2\sigma^2}\right),
    \quad x \in \R.
\]
When $\mu = 0$ and $\sigma^2 = 1$, we write $Z \sim \mathcal{N}(0,1)$ and call $Z$ a
\emph{standard normal} random variable.
\end{definition}

\begin{theorem}[The Gaussian Integral]
\label{thm:gaussian_integral}
$\displaystyle \int_{-\infty}^\infty e^{-x^2/2}\,dx = \sqrt{2\pi}$.
Consequently, the standard normal PDF integrates to $1$.
\end{theorem}

\begin{proof}
Let $I = \int_{-\infty}^\infty e^{-x^2/2}\,dx$. We show $I^2 = 2\pi$. Consider
\[
    I^2 = \left(\int_{-\infty}^\infty e^{-x^2/2}\,dx\right)
          \left(\int_{-\infty}^\infty e^{-y^2/2}\,dy\right)
        = \int_{-\infty}^\infty \int_{-\infty}^\infty e^{-(x^2+y^2)/2}\,dx\,dy.
\]
Switch to polar coordinates: $x = r\cos\theta$, $y = r\sin\theta$,
$dx\,dy = r\,dr\,d\theta$:
\[
    I^2 = \int_0^{2\pi}\int_0^\infty e^{-r^2/2}\,r\,dr\,d\theta
    = 2\pi \int_0^\infty r\,e^{-r^2/2}\,dr.
\]
Substituting $u = r^2/2$, $du = r\,dr$:
\[
    \int_0^\infty r\,e^{-r^2/2}\,dr = \int_0^\infty e^{-u}\,du = 1.
\]
Therefore $I^2 = 2\pi$, so $I = \sqrt{2\pi}$ (since $I > 0$). It follows that
$\int_{-\infty}^\infty \frac{1}{\sqrt{2\pi}}e^{-x^2/2}\,dx = 1$.
\end{proof}

\begin{proposition}
\label{prop:normal_moments}
If $X \sim \mathcal{N}(\mu, \sigma^2)$, then $\E[X] = \mu$ and $\Var(X) = \sigma^2$.
\end{proposition}

\begin{proof}
Let $Z = (X - \mu)/\sigma$, so $Z \sim \mathcal{N}(0,1)$ and $X = \mu + \sigma Z$.
It suffices to show $\E[Z] = 0$ and $\Var(Z) = 1$.

For the mean: $\E[Z] = \frac{1}{\sqrt{2\pi}}\int_{-\infty}^\infty z\,e^{-z^2/2}\,dz = 0$,
since the integrand $z \mapsto z\,e^{-z^2/2}$ is an odd function.

For the variance: $\E[Z^2] = \frac{1}{\sqrt{2\pi}}\int_{-\infty}^\infty z^2 e^{-z^2/2}\,dz$.
Integrate by parts with $u = z$, $dv = z\,e^{-z^2/2}\,dz$, so $du = dz$ and
$v = -e^{-z^2/2}$:
\[
    \int_{-\infty}^\infty z^2 e^{-z^2/2}\,dz
    = \left[-z\,e^{-z^2/2}\right]_{-\infty}^\infty + \int_{-\infty}^\infty e^{-z^2/2}\,dz
    = 0 + \sqrt{2\pi}.
\]
Thus $\E[Z^2] = 1$, and $\Var(Z) = \E[Z^2] - (\E[Z])^2 = 1 - 0 = 1$. Therefore
$\E[X] = \mu + \sigma \cdot 0 = \mu$ and $\Var(X) = \sigma^2 \Var(Z) = \sigma^2$.
\end{proof}

%% ============================================================
%% SECTION 6: EXPECTATION
%% ============================================================
\section{Expectation}
\label{sec:expectation}

\begin{definition}[Expectation]
\label{def:expectation}
Let $X$ be a random variable.
\begin{itemize}[nosep]
    \item \textbf{Discrete case:} If $X$ is discrete with PMF $p_X$, then
    \[
        \E[X] = \sum_x x \, p_X(x),
    \]
    provided the sum converges absolutely.
    \item \textbf{Continuous case:} If $X$ is continuous with PDF $f_X$, then
    \[
        \E[X] = \int_{-\infty}^\infty x \, f_X(x)\,dx,
    \]
    provided the integral converges absolutely.
\end{itemize}
\end{definition}

\begin{theorem}[Law of the Unconscious Statistician (LOTUS)]
\label{thm:lotus}
Let $g \colon \R \to \R$.
\begin{itemize}[nosep]
    \item \textbf{Discrete:} $\E[g(X)] = \sum_x g(x)\,p_X(x)$.
    \item \textbf{Continuous:} $\E[g(X)] = \int_{-\infty}^\infty g(x)\,f_X(x)\,dx$.
\end{itemize}
\end{theorem}

\begin{proof}[Proof (discrete case)]
Let $Y = g(X)$. For each value $y$ in the range of $g$,
$p_Y(y) = \Prob(Y = y) = \sum_{x : g(x) = y} p_X(x)$.
Then
\[
    \E[Y] = \sum_y y \, p_Y(y) = \sum_y y \sum_{x : g(x) = y} p_X(x)
    = \sum_y \sum_{x : g(x) = y} g(x)\,p_X(x) = \sum_x g(x)\,p_X(x),
\]
where the last step rearranges the double sum: every $x$ appears exactly once in the
partition $\{x : g(x) = y\}$ as $y$ ranges over the values of $g$.
\end{proof}

\begin{theorem}[Linearity of Expectation]
\label{thm:linearity}
For any random variables $X$ and $Y$ (on the same sample space) and constants $a, b \in \R$,
\[
    \E[aX + bY] = a\E[X] + b\E[Y].
\]
This holds \textbf{without any independence assumption}.
\end{theorem}

\begin{proof}
We prove the discrete case; the continuous case is analogous with integrals.

First, $\E[aX] = \sum_x (ax)\,p_X(x) = a\sum_x x\,p_X(x) = a\E[X]$.

For the sum: let $(X, Y)$ have joint PMF $p_{X,Y}$. Then
\begin{align*}
    \E[X + Y] &= \sum_x \sum_y (x + y)\,p_{X,Y}(x,y) \\
    &= \sum_x \sum_y x\,p_{X,Y}(x,y) + \sum_x \sum_y y\,p_{X,Y}(x,y) \\
    &= \sum_x x \sum_y p_{X,Y}(x,y) + \sum_y y \sum_x p_{X,Y}(x,y) \\
    &= \sum_x x\,p_X(x) + \sum_y y\,p_Y(y) = \E[X] + \E[Y].
\end{align*}
Combining: $\E[aX + bY] = \E[aX] + \E[bY] = a\E[X] + b\E[Y]$.
\end{proof}

\begin{proposition}[Expectation of Indicator Random Variables]
\label{prop:indicator_expectation}
Let $A$ be an event, and define $\indicator_A(\omega) = 1$ if $\omega \in A$ and
$\indicator_A(\omega) = 0$ otherwise. Then $\E[\indicator_A] = \Prob(A)$.
\end{proposition}

\begin{proof}
$\indicator_A \sim \mathrm{Bernoulli}(\Prob(A))$, so
$\E[\indicator_A] = 1 \cdot \Prob(A) + 0 \cdot (1 - \Prob(A)) = \Prob(A)$.
\end{proof}

\begin{keyresult}[Linearity of Expectation]
Linearity of expectation is arguably the single most useful property in probability. It
holds unconditionally --- no independence required. In RL, the value function
$V^\pi(s) = \E\!\left[\sum_{t=0}^\infty \gamma^t R_t \mid S_0 = s\right]$ is linear in the
reward, and linearity lets us decompose it step by step.
\end{keyresult}

%% ============================================================
%% SECTION 7: VARIANCE AND COVARIANCE
%% ============================================================
\section{Variance and Covariance}
\label{sec:variance}

\begin{definition}[Variance]
\label{def:variance}
The \emph{variance} of a random variable $X$ is
\[
    \Var(X) = \E\!\left[(X - \E[X])^2\right],
\]
provided this expectation exists.
\end{definition}

\begin{theorem}[Computational Formula for Variance]
\label{thm:var_formula}
$\Var(X) = \E[X^2] - (\E[X])^2$.
\end{theorem}

\begin{proof}
Let $\mu = \E[X]$. Expanding the square:
\[
    \Var(X) = \E[(X - \mu)^2] = \E[X^2 - 2\mu X + \mu^2]
    = \E[X^2] - 2\mu\E[X] + \mu^2 = \E[X^2] - 2\mu^2 + \mu^2 = \E[X^2] - \mu^2.
    \qedhere
\]
\end{proof}

\begin{proposition}[Affine Transformation of Variance]
\label{prop:var_affine}
For constants $a, b \in \R$, $\Var(aX + b) = a^2\Var(X)$.
\end{proposition}

\begin{proof}
$\E[aX + b] = a\E[X] + b$, so
\[
    \Var(aX + b) = \E[(aX + b - a\E[X] - b)^2] = \E[a^2(X - \E[X])^2]
    = a^2\E[(X - \E[X])^2] = a^2\Var(X).
    \qedhere
\]
\end{proof}

\begin{definition}[Standard Deviation]
\label{def:std_dev}
The \emph{standard deviation} of $X$ is $\mathrm{SD}(X) = \sqrt{\Var(X)}$.
\end{definition}

\begin{definition}[Covariance]
\label{def:covariance}
The \emph{covariance} of random variables $X$ and $Y$ is
\[
    \Cov(X, Y) = \E\!\left[(X - \E[X])(Y - \E[Y])\right].
\]
\end{definition}

\begin{theorem}[Computational Formula for Covariance]
\label{thm:cov_formula}
$\Cov(X, Y) = \E[XY] - \E[X]\E[Y]$.
\end{theorem}

\begin{proof}
Let $\mu_X = \E[X]$ and $\mu_Y = \E[Y]$. Expanding:
\begin{align*}
    \Cov(X,Y) &= \E[(X - \mu_X)(Y - \mu_Y)] = \E[XY - \mu_Y X - \mu_X Y + \mu_X\mu_Y] \\
    &= \E[XY] - \mu_Y\E[X] - \mu_X\E[Y] + \mu_X\mu_Y = \E[XY] - \mu_X\mu_Y.
    \qedhere
\end{align*}
\end{proof}

\begin{proposition}[Properties of Covariance]
\label{prop:cov_properties}
For random variables $X, Y, Z$ and constants $a, b \in \R$:
\begin{enumerate}[nosep]
    \item $\Cov(X, X) = \Var(X)$.
    \item $\Cov(X, Y) = \Cov(Y, X)$ \hfill(symmetry).
    \item $\Cov(aX + bY, Z) = a\Cov(X, Z) + b\Cov(Y, Z)$ \hfill(bilinearity).
\end{enumerate}
\end{proposition}

\begin{proof}
\textbf{(1)} $\Cov(X,X) = \E[(X - \E[X])(X - \E[X])] = \E[(X - \E[X])^2] = \Var(X)$.

\textbf{(2)} Immediate from $\E[XY] = \E[YX]$.

\textbf{(3)} Using the computational formula:
\begin{align*}
    \Cov(aX + bY, Z) &= \E[(aX + bY)Z] - \E[aX + bY]\E[Z] \\
    &= a\E[XZ] + b\E[YZ] - (a\E[X] + b\E[Y])\E[Z] \\
    &= a(\E[XZ] - \E[X]\E[Z]) + b(\E[YZ] - \E[Y]\E[Z]) \\
    &= a\Cov(X,Z) + b\Cov(Y,Z).
    \qedhere
\end{align*}
\end{proof}

\begin{theorem}[Variance of a Sum]
\label{thm:var_sum}
$\Var(X + Y) = \Var(X) + \Var(Y) + 2\Cov(X,Y)$.
\end{theorem}

\begin{proof}
\begin{align*}
    \Var(X + Y) &= \Cov(X + Y, X + Y) \\
    &= \Cov(X, X) + \Cov(X, Y) + \Cov(Y, X) + \Cov(Y, Y) \\
    &= \Var(X) + 2\Cov(X, Y) + \Var(Y),
\end{align*}
where the second line uses bilinearity and the third uses
$\Cov(X,X) = \Var(X)$, $\Cov(Y,Y) = \Var(Y)$, and symmetry.
\end{proof}

\begin{theorem}[Independent Implies Uncorrelated]
\label{thm:indep_uncorr}
If $X$ and $Y$ are independent random variables, then $\Cov(X,Y) = 0$.
\end{theorem}

\begin{proof}
For independent discrete random variables,
\[
    \E[XY] = \sum_x \sum_y xy\,p_{X,Y}(x,y) = \sum_x \sum_y xy\,p_X(x)\,p_Y(y)
    = \left(\sum_x x\,p_X(x)\right)\!\left(\sum_y y\,p_Y(y)\right) = \E[X]\E[Y].
\]
By Theorem~\ref{thm:cov_formula}, $\Cov(X,Y) = \E[XY] - \E[X]\E[Y] = 0$.
\end{proof}

\begin{corollary}[Variance of Sum of Independent Random Variables]
\label{cor:var_indep_sum}
If $X_1, \ldots, X_n$ are independent, then $\Var(X_1 + \cdots + X_n) = \Var(X_1) + \cdots + \Var(X_n)$.
\end{corollary}

\begin{proof}
By induction using Theorem~\ref{thm:var_sum} and Theorem~\ref{thm:indep_uncorr}: all
cross-covariance terms vanish.
\end{proof}

\begin{definition}[Correlation Coefficient]
\label{def:correlation}
The \emph{correlation} of $X$ and $Y$ (with $\Var(X), \Var(Y) > 0$) is
\[
    \rho(X,Y) = \frac{\Cov(X,Y)}{\sqrt{\Var(X)\Var(Y)}}.
\]
One can show $-1 \leq \rho(X,Y) \leq 1$ (this follows from the Cauchy--Schwarz inequality,
Theorem~\ref{thm:cauchy_schwarz}).
\end{definition}

\subsection{Variances of the Standard Distributions}

We collect the variance computations for the distributions of Sections~\ref{sec:discrete_rv}
and~\ref{sec:continuous_rv}. The proofs of $\Var(X) = p(1-p)$ for Bernoulli,
$\Var(X) = np(1-p)$ for Binomial, $\Var(X) = (1-p)/p^2$ for Geometric, and
$\Var(X) = \lambda$ for Poisson were given in those sections. Similarly,
$\Var(X) = (b-a)^2/12$ for Uniform, $\Var(X) = 1/\lambda^2$ for Exponential, and
$\Var(X) = \sigma^2$ for Normal were proved there.

\begin{center}
\begin{tabular}{lcc}
\toprule
Distribution & $\E[X]$ & $\Var(X)$ \\
\midrule
$\mathrm{Bernoulli}(p)$ & $p$ & $p(1-p)$ \\
$\mathrm{Binomial}(n,p)$ & $np$ & $np(1-p)$ \\
$\mathrm{Geometric}(p)$ & $1/p$ & $(1-p)/p^2$ \\
$\mathrm{Poisson}(\lambda)$ & $\lambda$ & $\lambda$ \\
$\mathrm{Uniform}(a,b)$ & $(a+b)/2$ & $(b-a)^2/12$ \\
$\mathrm{Exponential}(\lambda)$ & $1/\lambda$ & $1/\lambda^2$ \\
$\mathcal{N}(\mu, \sigma^2)$ & $\mu$ & $\sigma^2$ \\
\bottomrule
\end{tabular}
\end{center}

%% ============================================================
%% SECTION 8: JOINT DISTRIBUTIONS
%% ============================================================
\section{Joint Distributions}
\label{sec:joint}

\begin{definition}[Joint PMF]
\label{def:joint_pmf}
Discrete random variables $X$ and $Y$ have \emph{joint PMF}
\[
    p_{X,Y}(x,y) = \Prob(X = x, Y = y).
\]
\end{definition}

\begin{definition}[Joint PDF]
\label{def:joint_pdf}
Continuous random variables $X$ and $Y$ have \emph{joint PDF} $f_{X,Y}$ if
\[
    \Prob((X,Y) \in A) = \iint_A f_{X,Y}(x,y)\,dx\,dy
\]
for every ``reasonable'' region $A \subseteq \R^2$.
\end{definition}

\subsection{Marginal Distributions}

\begin{proposition}[Marginal from Joint]
\label{prop:marginals}
\begin{itemize}[nosep]
    \item \textbf{Discrete:} $p_X(x) = \sum_y p_{X,Y}(x,y)$.
    \item \textbf{Continuous:} $f_X(x) = \int_{-\infty}^\infty f_{X,Y}(x,y)\,dy$.
\end{itemize}
\end{proposition}

\begin{proof}
In the discrete case, $p_X(x) = \Prob(X = x) = \Prob\!\left(\bigcup_y \{X = x, Y = y\}\right)
= \sum_y \Prob(X = x, Y = y) = \sum_y p_{X,Y}(x,y)$, where the third equality uses
countable additivity (the events $\{Y = y\}$ are disjoint). The continuous case is analogous
with integration replacing summation.
\end{proof}

\subsection{Conditional Distributions}

\begin{definition}[Conditional PMF]
\label{def:cond_pmf}
If $\Prob(X = x) > 0$, the \emph{conditional PMF of $Y$ given $X = x$} is
\[
    p_{Y|X}(y \mid x) = \Prob(Y = y \mid X = x) = \frac{p_{X,Y}(x,y)}{p_X(x)}.
\]
\end{definition}

\begin{definition}[Conditional PDF]
\label{def:cond_pdf}
If $f_X(x) > 0$, the \emph{conditional PDF of $Y$ given $X = x$} is
\[
    f_{Y|X}(y \mid x) = \frac{f_{X,Y}(x,y)}{f_X(x)}.
\]
\end{definition}

\subsection{Independence of Random Variables}

\begin{definition}[Independence of Random Variables]
\label{def:indep_rv}
Random variables $X$ and $Y$ are \emph{independent} if their joint distribution factors:
\begin{itemize}[nosep]
    \item \textbf{Discrete:} $p_{X,Y}(x,y) = p_X(x)\,p_Y(y)$ for all $x, y$.
    \item \textbf{Continuous:} $f_{X,Y}(x,y) = f_X(x)\,f_Y(y)$ for all $x, y$.
\end{itemize}
\end{definition}

\subsection{Conditional Expectation}

\begin{definition}[Conditional Expectation Given an Event]
\label{def:cond_exp}
In the discrete case, the \emph{conditional expectation of $Y$ given $X = x$} is
\[
    \E[Y \mid X = x] = \sum_y y \, p_{Y|X}(y \mid x).
\]
In the continuous case,
\[
    \E[Y \mid X = x] = \int_{-\infty}^\infty y \, f_{Y|X}(y \mid x)\,dy.
\]
The function $g(x) = \E[Y \mid X = x]$ defines a new random variable $\E[Y \mid X] = g(X)$.
\end{definition}

\begin{theorem}[Law of Total Expectation (Tower Property)]
\label{thm:total_expectation}
$\E[Y] = \E\!\left[\E[Y \mid X]\right]$.
\end{theorem}

\begin{proof}[Proof (discrete case)]
Let $g(x) = \E[Y \mid X = x]$. Then
\begin{align*}
    \E[\E[Y \mid X]] &= \E[g(X)] = \sum_x g(x)\,p_X(x)
    = \sum_x \left(\sum_y y\,p_{Y|X}(y \mid x)\right) p_X(x) \\
    &= \sum_x \sum_y y \cdot \frac{p_{X,Y}(x,y)}{p_X(x)} \cdot p_X(x)
    = \sum_x \sum_y y\,p_{X,Y}(x,y) \\
    &= \sum_y y \sum_x p_{X,Y}(x,y) = \sum_y y\,p_Y(y) = \E[Y].
    \qedhere
\end{align*}
\end{proof}

\begin{theorem}[Law of Total Variance]
\label{thm:total_variance}
$\Var(Y) = \E[\Var(Y \mid X)] + \Var(\E[Y \mid X])$.
\end{theorem}

\begin{proof}
We use $\Var(Y) = \E[Y^2] - (\E[Y])^2$.

The conditional variance is $\Var(Y \mid X) = \E[Y^2 \mid X] - (\E[Y \mid X])^2$.
Taking expectations of both sides and using the tower property:
\[
    \E[\Var(Y \mid X)] = \E[\E[Y^2 \mid X]] - \E[(\E[Y \mid X])^2]
    = \E[Y^2] - \E[(\E[Y \mid X])^2].
\]
For the second term, $\Var(\E[Y \mid X]) = \E[(\E[Y \mid X])^2] - (\E[\E[Y \mid X]])^2
= \E[(\E[Y \mid X])^2] - (\E[Y])^2$.
Adding:
\begin{align*}
    \E[\Var(Y \mid X)] + \Var(\E[Y \mid X])
    &= \left(\E[Y^2] - \E[(\E[Y \mid X])^2]\right) + \left(\E[(\E[Y \mid X])^2] - (\E[Y])^2\right) \\
    &= \E[Y^2] - (\E[Y])^2 = \Var(Y).
    \qedhere
\end{align*}
\end{proof}

\begin{keyresult}[Tower Property and Dynamic Programming]
The law of total expectation (tower property) is the engine of dynamic programming. The
Bellman equation
\[
    V(s) = \E[R + \gamma V(S') \mid S = s]
    = \sum_a \pi(a \mid s) \sum_{s'} P(s' \mid s, a)\left[R(s,a,s') + \gamma V(s')\right]
\]
is obtained by conditioning on the action and next state. Phase~01 will prove this
rigorously with $\sigma$-algebras.
\end{keyresult}

%% ============================================================
%% SECTION 9: MOMENT GENERATING FUNCTIONS
%% ============================================================
\section{Moment Generating Functions}
\label{sec:mgf}

\begin{definition}[Moment Generating Function]
\label{def:mgf}
The \emph{moment generating function (MGF)} of a random variable $X$ is
\[
    M_X(t) = \E[e^{tX}],
\]
defined for all $t$ in an open interval containing $0$ where this expectation is finite.
\end{definition}

\begin{theorem}[MGF Determines the Distribution]
\label{thm:mgf_unique}
If $M_X(t) = M_Y(t)$ for all $t$ in an open interval containing $0$, then $X$ and $Y$
have the same distribution.
\end{theorem}

\begin{remark}
The proof uses the theory of Laplace transforms and analytic continuation. We state this
theorem without proof and use it as a tool for identifying distributions.
\end{remark}

\begin{theorem}[Moments from MGF]
\label{thm:mgf_moments}
If $M_X(t)$ exists in a neighborhood of $0$, then
\[
    \E[X^n] = M_X^{(n)}(0) = \frac{d^n}{dt^n} M_X(t)\bigg|_{t=0}
\]
for every $n \geq 1$.
\end{theorem}

\begin{proof}
Since $e^{tX} = \sum_{k=0}^\infty \frac{(tX)^k}{k!}$ and the MGF exists in a neighborhood
of $0$ (justifying interchange of expectation and summation), we have
\[
    M_X(t) = \E[e^{tX}] = \sum_{k=0}^\infty \frac{t^k}{k!}\E[X^k].
\]
Differentiating $n$ times with respect to $t$:
\[
    M_X^{(n)}(t) = \sum_{k=n}^\infty \frac{t^{k-n}}{(k-n)!}\E[X^k].
\]
Setting $t = 0$: $M_X^{(n)}(0) = \E[X^n]$ (only the $k = n$ term survives).
\end{proof}

\begin{theorem}[MGF of Sum of Independent Random Variables]
\label{thm:mgf_sum}
If $X$ and $Y$ are independent, then $M_{X+Y}(t) = M_X(t) \cdot M_Y(t)$.
\end{theorem}

\begin{proof}
$M_{X+Y}(t) = \E[e^{t(X+Y)}] = \E[e^{tX} \cdot e^{tY}] = \E[e^{tX}]\cdot\E[e^{tY}]
= M_X(t) \cdot M_Y(t)$,
where the third equality uses the fact that for independent random variables,
$\E[g(X)h(Y)] = \E[g(X)]\E[h(Y)]$ (with $g(x) = e^{tx}$ and $h(y) = e^{ty}$).
\end{proof}

\subsection{MGFs of Standard Distributions}

\begin{proposition}[MGF of Bernoulli]
\label{prop:mgf_bernoulli}
If $X \sim \mathrm{Bernoulli}(p)$, then $M_X(t) = 1 - p + pe^t$.
\end{proposition}

\begin{proof}
$M_X(t) = \E[e^{tX}] = e^{t \cdot 0}(1-p) + e^{t \cdot 1} p = (1-p) + pe^t$.
\end{proof}

\begin{proposition}[MGF of Poisson]
\label{prop:mgf_poisson}
If $X \sim \mathrm{Poisson}(\lambda)$, then $M_X(t) = \exp(\lambda(e^t - 1))$.
\end{proposition}

\begin{proof}
\[
    M_X(t) = \sum_{k=0}^\infty e^{tk}\frac{e^{-\lambda}\lambda^k}{k!}
    = e^{-\lambda}\sum_{k=0}^\infty \frac{(\lambda e^t)^k}{k!}
    = e^{-\lambda}\cdot e^{\lambda e^t} = e^{\lambda(e^t - 1)}.
    \qedhere
\]
\end{proof}

\begin{proposition}[MGF of Normal]
\label{prop:mgf_normal}
If $X \sim \mathcal{N}(\mu, \sigma^2)$, then $M_X(t) = \exp(\mu t + \sigma^2 t^2/2)$.
\end{proposition}

\begin{proof}
It suffices to prove the case $Z \sim \mathcal{N}(0,1)$, since
$M_X(t) = \E[e^{t(\mu + \sigma Z)}] = e^{\mu t} M_Z(\sigma t)$.

For $Z \sim \mathcal{N}(0,1)$:
\begin{align*}
    M_Z(t) &= \frac{1}{\sqrt{2\pi}}\int_{-\infty}^\infty e^{tz} e^{-z^2/2}\,dz
    = \frac{1}{\sqrt{2\pi}}\int_{-\infty}^\infty e^{-(z^2 - 2tz)/2}\,dz \\
    &= \frac{1}{\sqrt{2\pi}}\int_{-\infty}^\infty e^{-((z-t)^2 - t^2)/2}\,dz
    = e^{t^2/2} \cdot \frac{1}{\sqrt{2\pi}}\int_{-\infty}^\infty e^{-(z-t)^2/2}\,dz
    = e^{t^2/2},
\end{align*}
where the last integral equals $1$ by the substitution $u = z - t$ and
Theorem~\ref{thm:gaussian_integral}. Thus $M_X(t) = e^{\mu t} \cdot e^{(\sigma t)^2/2}
= e^{\mu t + \sigma^2 t^2/2}$.
\end{proof}

\subsection{Applications of MGFs}

\begin{theorem}[Sum of Independent Poissons]
\label{thm:sum_poisson}
If $X \sim \mathrm{Poisson}(\lambda)$ and $Y \sim \mathrm{Poisson}(\mu)$ are independent,
then $X + Y \sim \mathrm{Poisson}(\lambda + \mu)$.
\end{theorem}

\begin{proof}
$M_{X+Y}(t) = M_X(t)M_Y(t) = e^{\lambda(e^t-1)} \cdot e^{\mu(e^t-1)}
= e^{(\lambda+\mu)(e^t-1)}$,
which is the MGF of $\mathrm{Poisson}(\lambda + \mu)$. By
Theorem~\ref{thm:mgf_unique}, $X + Y \sim \mathrm{Poisson}(\lambda + \mu)$.
\end{proof}

\begin{theorem}[Sum of Independent Normals]
\label{thm:sum_normal}
If $X \sim \mathcal{N}(\mu_1, \sigma_1^2)$ and $Y \sim \mathcal{N}(\mu_2, \sigma_2^2)$
are independent, then $X + Y \sim \mathcal{N}(\mu_1 + \mu_2, \sigma_1^2 + \sigma_2^2)$.
\end{theorem}

\begin{proof}
$M_{X+Y}(t) = M_X(t)M_Y(t) = e^{\mu_1 t + \sigma_1^2 t^2/2} \cdot
e^{\mu_2 t + \sigma_2^2 t^2/2} = e^{(\mu_1+\mu_2)t + (\sigma_1^2+\sigma_2^2)t^2/2}$,
which is the MGF of $\mathcal{N}(\mu_1 + \mu_2, \sigma_1^2 + \sigma_2^2)$. The result
follows from Theorem~\ref{thm:mgf_unique}.
\end{proof}

%% ============================================================
%% SECTION 10: PROBABILITY INEQUALITIES
%% ============================================================
\section{Probability Inequalities}
\label{sec:inequalities}

\begin{theorem}[Markov's Inequality]
\label{thm:markov}
If $X \geq 0$ and $a > 0$, then
\[
    \Prob(X \geq a) \leq \frac{\E[X]}{a}.
\]
\end{theorem}

\begin{proof}
We have $X \geq a \cdot \indicator_{\{X \geq a\}}$, since: if $X \geq a$ then the right
side is $a \leq X$, and if $X < a$ then the right side is $0 \leq X$. Taking expectations:
\[
    \E[X] \geq a \cdot \E[\indicator_{\{X \geq a\}}] = a \cdot \Prob(X \geq a).
\]
Dividing by $a > 0$ gives the result.
\end{proof}

\begin{theorem}[Chebyshev's Inequality]
\label{thm:chebyshev}
For any random variable $X$ with finite variance and any $k > 0$,
\[
    \Prob(|X - \E[X]| \geq k) \leq \frac{\Var(X)}{k^2}.
\]
\end{theorem}

\begin{proof}
Apply Markov's inequality to the non-negative random variable $(X - \E[X])^2$ with
threshold $k^2$:
\[
    \Prob(|X - \E[X]| \geq k) = \Prob((X - \E[X])^2 \geq k^2)
    \leq \frac{\E[(X - \E[X])^2]}{k^2} = \frac{\Var(X)}{k^2}.
    \qedhere
\]
\end{proof}

\begin{theorem}[Cauchy--Schwarz Inequality for Expectations]
\label{thm:cauchy_schwarz}
For any random variables $X$ and $Y$ with $\E[X^2], \E[Y^2] < \infty$,
\[
    |\E[XY]|^2 \leq \E[X^2]\,\E[Y^2].
\]
\end{theorem}

\begin{proof}
If $\E[Y^2] = 0$, then $Y = 0$ a.s., so $\E[XY] = 0$ and both sides are $0$.

Assume $\E[Y^2] > 0$. For any $t \in \R$, $\E[(X + tY)^2] \geq 0$. Expanding:
\[
    \E[X^2] + 2t\,\E[XY] + t^2\,\E[Y^2] \geq 0 \quad \text{for all } t \in \R.
\]
This is a quadratic in $t$ that is non-negative for all $t$. A quadratic $at^2 + bt + c$
with $a > 0$ is non-negative for all $t$ if and only if its discriminant satisfies
$b^2 - 4ac \leq 0$. Here $a = \E[Y^2] > 0$, $b = 2\E[XY]$, $c = \E[X^2]$, so
\[
    4(\E[XY])^2 - 4\E[X^2]\E[Y^2] \leq 0,
\]
which gives $(\E[XY])^2 \leq \E[X^2]\E[Y^2]$.
\end{proof}

\begin{theorem}[Jensen's Inequality]
\label{thm:jensen}
If $g \colon \R \to \R$ is convex and $\E[|X|] < \infty$, then
\[
    g(\E[X]) \leq \E[g(X)].
\]
\end{theorem}

\begin{proof}[Proof (discrete case, by induction)]
We prove that if $X$ takes values $x_1, \ldots, x_n$ with probabilities $p_1, \ldots, p_n$
(where $p_i > 0$ and $\sum p_i = 1$), then $g\!\left(\sum_{i=1}^n p_i x_i\right) \leq
\sum_{i=1}^n p_i\,g(x_i)$.

\emph{Base case ($n = 1$):} $g(p_1 x_1) = g(x_1) = p_1 g(x_1)$. Holds trivially.

\emph{Base case ($n = 2$):} This is the definition of convexity:
$g(p_1 x_1 + p_2 x_2) \leq p_1 g(x_1) + p_2 g(x_2)$ with $p_1 + p_2 = 1$.

\emph{Inductive step:} Assume the result for $n - 1$ values. Let $q = 1 - p_n > 0$
(the case $p_n = 1$ is trivial). Write
\[
    \sum_{i=1}^n p_i x_i = q \cdot \underbrace{\sum_{i=1}^{n-1}\frac{p_i}{q}x_i}_{= \bar{x}} + p_n x_n.
\]
Since $q + p_n = 1$, convexity gives
\[
    g\!\left(\sum_{i=1}^n p_i x_i\right) = g(q\bar{x} + p_n x_n)
    \leq q\,g(\bar{x}) + p_n\,g(x_n).
\]
The weights $p_i/q$ for $i = 1, \ldots, n-1$ sum to $1$, so the inductive hypothesis
gives $g(\bar{x}) \leq \sum_{i=1}^{n-1}\frac{p_i}{q}g(x_i)$. Substituting:
\[
    g\!\left(\sum_{i=1}^n p_i x_i\right)
    \leq q \sum_{i=1}^{n-1}\frac{p_i}{q}g(x_i) + p_n\,g(x_n)
    = \sum_{i=1}^n p_i\,g(x_i).
    \qedhere
\]
\end{proof}

\begin{rlconnection}[Concentration Inequalities in RL]
Concentration inequalities are the quantitative tools that turn convergence concepts into
convergence rates. Markov and Chebyshev give basic tail bounds. Phase~01 will add the
sharper Hoeffding and Azuma--Hoeffding inequalities used in RL convergence proofs.
\end{rlconnection}

%% ============================================================
%% SECTION 11: LIMIT THEOREMS
%% ============================================================
\section{Limit Theorems}
\label{sec:limit_theorems}

\begin{definition}[Convergence in Probability]
\label{def:conv_prob}
A sequence of random variables $X_1, X_2, \ldots$ \emph{converges in probability} to a
random variable $X$, written $X_n \xrightarrow{\Prob} X$, if for every $\varepsilon > 0$,
\[
    \lim_{n \to \infty} \Prob(|X_n - X| \geq \varepsilon) = 0.
\]
\end{definition}

\begin{theorem}[Weak Law of Large Numbers (WLLN)]
\label{thm:wlln}
Let $X_1, X_2, \ldots$ be independent and identically distributed (i.i.d.) random variables
with $\E[X_i] = \mu$ and $\Var(X_i) = \sigma^2 < \infty$. Let
$\bar{X}_n = \frac{1}{n}\sum_{i=1}^n X_i$. Then $\bar{X}_n \xrightarrow{\Prob} \mu$.
\end{theorem}

\begin{proof}
By linearity of expectation, $\E[\bar{X}_n] = \mu$. By independence,
\[
    \Var(\bar{X}_n) = \Var\!\left(\frac{1}{n}\sum_{i=1}^n X_i\right)
    = \frac{1}{n^2}\sum_{i=1}^n \Var(X_i) = \frac{\sigma^2}{n}.
\]
By Chebyshev's inequality (Theorem~\ref{thm:chebyshev}), for any $\varepsilon > 0$:
\[
    \Prob(|\bar{X}_n - \mu| \geq \varepsilon)
    \leq \frac{\Var(\bar{X}_n)}{\varepsilon^2} = \frac{\sigma^2}{n\varepsilon^2}
    \to 0 \quad \text{as } n \to \infty.
    \qedhere
\]
\end{proof}

\begin{theorem}[Strong Law of Large Numbers (SLLN)]
\label{thm:slln}
Under the same hypotheses as Theorem~\ref{thm:wlln}, the convergence holds almost surely:
$\bar{X}_n \to \mu$ with probability $1$.
\end{theorem}

\begin{remark}
The SLLN requires measure-theoretic tools (the Borel--Cantelli lemma and
truncation arguments) and is proved in Phase~01.
\end{remark}

\begin{definition}[Convergence in Distribution]
\label{def:conv_dist}
A sequence $X_1, X_2, \ldots$ \emph{converges in distribution} to $X$, written
$X_n \xrightarrow{d} X$, if
\[
    \lim_{n \to \infty} F_{X_n}(x) = F_X(x)
\]
at every point $x$ where $F_X$ is continuous.
\end{definition}

\begin{theorem}[Central Limit Theorem (CLT)]
\label{thm:clt}
Let $X_1, X_2, \ldots$ be i.i.d.\ with $\E[X_i] = \mu$ and $\Var(X_i) = \sigma^2 \in
(0, \infty)$. Then
\[
    \frac{\bar{X}_n - \mu}{\sigma/\sqrt{n}} = \frac{\sum_{i=1}^n X_i - n\mu}{\sigma\sqrt{n}}
    \xrightarrow{d} \mathcal{N}(0, 1).
\]
\end{theorem}

\begin{remark}
The standard proof uses characteristic functions (or equivalently, MGFs when they exist)
and is deferred to Phase~01. Here we state the result and illustrate its use.
\end{remark}

\begin{example}[CLT Application: Polling]
\label{ex:clt_polling}
A poll surveys $n = 1000$ voters, each independently supporting a candidate with unknown
probability $p$. The sample proportion is $\hat{p} = \bar{X}_n$ where $X_i \sim
\mathrm{Bernoulli}(p)$. By the CLT,
\[
    \frac{\hat{p} - p}{\sqrt{p(1-p)/n}} \approx \mathcal{N}(0, 1).
\]
A $95\%$ confidence interval for $p$ is approximately $\hat{p} \pm 1.96\sqrt{\hat{p}(1-\hat{p})/n}$.
With $n = 1000$ and $\hat{p} = 0.52$, the margin of error is
$1.96\sqrt{0.52 \cdot 0.48/1000} \approx 0.031$.
\end{example}

\begin{example}[CLT Application: Sum of Dice]
\label{ex:clt_dice}
Roll a fair die $n = 100$ times and let $S_n = X_1 + \cdots + X_{100}$ be the total.
Each $X_i$ is uniform on $\{1, \ldots, 6\}$ with $\mu = 7/2$ and
$\sigma^2 = 35/12$. By the CLT,
\[
    \frac{S_{100} - 350}{\sqrt{100 \cdot 35/12}} \approx \mathcal{N}(0,1).
\]
So $\Prob(S_{100} \leq 330) \approx \Phi\!\left(\frac{330 - 350}{\sqrt{2916.67/100}}\right)
= \Phi(-1.17) \approx 0.121$.
\end{example}

\begin{keyresult}[LLN and Monte Carlo]
The law of large numbers justifies Monte Carlo estimation: the sample average of rewards
converges to the expected return. This is why Monte Carlo policy evaluation works. The CLT
tells us how fast: the error is approximately normal with standard deviation proportional to
$1/\sqrt{n}$, giving us confidence intervals for value estimates.
\end{keyresult}

%% ============================================================
%% SECTION 12: TAIL BOUNDS AND CONFIDENCE INTERVALS
%% ============================================================
\section{Tail Bounds and Confidence Intervals}
\label{sec:tail_bounds}

Chebyshev's inequality gives tail bounds in terms of the variance, but these are often loose.
For bounded random variables, we can do much better.

\subsection{Chebyshev Confidence Intervals}

\begin{proposition}[Chebyshev Confidence Interval]
\label{prop:chebyshev_ci}
Let $X_1, \ldots, X_n$ be i.i.d.\ with mean $\mu$ and variance $\sigma^2 < \infty$. For
any $\delta \in (0, 1)$,
\[
    \Prob\!\left(|\bar{X}_n - \mu| < \frac{\sigma}{\sqrt{n\delta}}\right) \geq 1 - \delta.
\]
\end{proposition}

\begin{proof}
By Chebyshev's inequality with $k = \sigma/\sqrt{n\delta}$:
\[
    \Prob\!\left(|\bar{X}_n - \mu| \geq \frac{\sigma}{\sqrt{n\delta}}\right)
    \leq \frac{\Var(\bar{X}_n)}{k^2} = \frac{\sigma^2/n}{\sigma^2/(n\delta)} = \delta.
\]
Taking complements gives the result.
\end{proof}

\subsection{Hoeffding's Inequality}

\begin{lemma}[Hoeffding's Lemma]
\label{lem:hoeffding}
If $X$ is a random variable with $\E[X] = 0$ and $a \leq X \leq b$ almost surely, then
for all $t > 0$,
\[
    \E[e^{tX}] \leq \exp\!\left(\frac{t^2(b-a)^2}{8}\right).
\]
\end{lemma}

\begin{proof}
Since $a \leq X \leq b$, we can write $X = \alpha b + (1 - \alpha)a$ where
$\alpha = \frac{X - a}{b - a} \in [0, 1]$. By convexity of $e^{tx}$:
\[
    e^{tX} \leq \frac{X - a}{b - a}e^{tb} + \frac{b - X}{b - a}e^{ta}.
\]
Taking expectations and using $\E[X] = 0$:
\[
    \E[e^{tX}] \leq \frac{-a}{b-a}e^{tb} + \frac{b}{b-a}e^{ta}.
\]
Let $p = \frac{-a}{b-a}$, $h = t(b-a)$. Then
\[
    \E[e^{tX}] \leq p\,e^{(1-p)h} + (1-p)\,e^{-ph} = e^{-ph}\left(pe^h + 1 - p\right).
\]
Define $\varphi(h) = -ph + \ln(pe^h + 1 - p)$. Then $\varphi(0) = 0$ and $\varphi'(0) = 0$
(verify: $\varphi'(h) = -p + \frac{pe^h}{pe^h + 1-p}$, and at $h=0$ this gives
$-p + p = 0$). Furthermore,
\[
    \varphi''(h) = \frac{pe^h(1-p)}{(pe^h + 1-p)^2} \leq \frac{1}{4},
\]
since $\frac{u(1-u)}{1} \leq 1/4$ for $u \in [0,1]$, applied to $u = \frac{pe^h}{pe^h+1-p}$.
By Taylor's theorem, $\varphi(h) = \varphi(0) + \varphi'(0)h + \frac{1}{2}\varphi''(c)h^2
\leq h^2/8$ for some $c$. Therefore $\E[e^{tX}] \leq e^{h^2/8} = e^{t^2(b-a)^2/8}$.
\end{proof}

\begin{theorem}[Hoeffding's Inequality]
\label{thm:hoeffding}
Let $X_1, \ldots, X_n$ be independent random variables with $a_i \leq X_i \leq b_i$
almost surely. Let $\bar{X}_n = \frac{1}{n}\sum_{i=1}^n X_i$. Then for any $\varepsilon > 0$,
\[
    \Prob\!\left(|\bar{X}_n - \E[\bar{X}_n]| \geq \varepsilon\right)
    \leq 2\exp\!\left(-\frac{2n^2\varepsilon^2}{\sum_{i=1}^n (b_i - a_i)^2}\right).
\]
\end{theorem}

\begin{proof}
Let $S_n = \sum_{i=1}^n X_i$ and $\mu = \E[S_n]$. We bound $\Prob(S_n - \mu \geq n\varepsilon)$;
the lower tail is symmetric.

For any $t > 0$, by Markov's inequality applied to $e^{t(S_n - \mu)}$:
\[
    \Prob(S_n - \mu \geq n\varepsilon)
    = \Prob(e^{t(S_n - \mu)} \geq e^{tn\varepsilon})
    \leq \frac{\E[e^{t(S_n - \mu)}]}{e^{tn\varepsilon}}.
\]
By independence, $\E[e^{t(S_n - \mu)}] = \prod_{i=1}^n \E[e^{t(X_i - \E[X_i])}]$. Each
centered variable $X_i - \E[X_i]$ satisfies $a_i - \E[X_i] \leq X_i - \E[X_i] \leq
b_i - \E[X_i]$, a range of $b_i - a_i$. By Hoeffding's Lemma,
$\E[e^{t(X_i - \E[X_i])}] \leq e^{t^2(b_i - a_i)^2/8}$. Therefore
\[
    \Prob(S_n - \mu \geq n\varepsilon) \leq \exp\!\left(-tn\varepsilon +
    \frac{t^2}{8}\sum_{i=1}^n(b_i - a_i)^2\right).
\]
Minimizing over $t$: the exponent is minimized at
$t = \frac{4n\varepsilon}{\sum_{i=1}^n(b_i-a_i)^2}$, giving
\[
    \Prob(S_n - \mu \geq n\varepsilon) \leq \exp\!\left(-\frac{2n^2\varepsilon^2}{\sum_{i=1}^n(b_i-a_i)^2}\right).
\]
The same bound holds for $\Prob(S_n - \mu \leq -n\varepsilon)$ by applying the argument to
$-X_i$. A union bound gives the two-sided result.
\end{proof}

\begin{corollary}[Hoeffding for i.i.d.\ Bounded Variables]
\label{cor:hoeffding_iid}
If $X_1, \ldots, X_n$ are i.i.d.\ with $a \leq X_i \leq b$, then
\[
    \Prob(|\bar{X}_n - \mu| \geq \varepsilon) \leq 2e^{-2n\varepsilon^2/(b-a)^2}.
\]
\end{corollary}

\begin{proof}
Set $a_i = a$ and $b_i = b$ for all $i$ in Theorem~\ref{thm:hoeffding}. Then
$\sum_{i=1}^n(b_i - a_i)^2 = n(b-a)^2$.
\end{proof}

\subsection{Sample Complexity}

\begin{proposition}[Sample Complexity from Hoeffding]
\label{prop:sample_complexity}
Let $X_1, \ldots, X_n$ be i.i.d.\ with $a \leq X_i \leq b$. To guarantee
$\Prob(|\bar{X}_n - \mu| \geq \varepsilon) \leq \delta$, it suffices to take
\[
    n \geq \frac{(b-a)^2}{2\varepsilon^2}\ln\frac{2}{\delta}.
\]
\end{proposition}

\begin{proof}
By Corollary~\ref{cor:hoeffding_iid}, we need $2e^{-2n\varepsilon^2/(b-a)^2} \leq \delta$.
This is equivalent to $e^{-2n\varepsilon^2/(b-a)^2} \leq \delta/2$, i.e.,
$-2n\varepsilon^2/(b-a)^2 \leq \ln(\delta/2) = -\ln(2/\delta)$.
Rearranging: $n \geq \frac{(b-a)^2}{2\varepsilon^2}\ln\frac{2}{\delta}$.
\end{proof}

\begin{remark}
The dependence $n = O(\varepsilon^{-2}\ln\delta^{-1})$ is the standard sample complexity
for Monte Carlo estimation of a bounded mean. The logarithmic dependence on $1/\delta$ is
much better than Chebyshev's $O(\varepsilon^{-2}\delta^{-1})$ bound.
\end{remark}

\begin{rlconnection}[Sample Complexity in RL]
In RL, if rewards are bounded in $[0, 1]$, the sample complexity of Monte Carlo policy
evaluation with accuracy $\varepsilon$ and confidence $1 - \delta$ is
$n \geq \frac{1}{2\varepsilon^2}\ln\frac{2}{\delta}$.
For $\varepsilon = 0.01$ and $\delta = 0.05$, this gives $n \geq 18{,}445$ samples.
Hoeffding's inequality makes this precise; Chebyshev would require
$n \geq 1/(\varepsilon^2 \delta) = 200{,}000$ --- an order of magnitude worse.
\end{rlconnection}

\end{document}
